
\section{Quantification \& Régression (concentrations, fractions de phases, dopage) (Section 9)}

\subsection{Objectif}
Estimer des quantités physico-chimiques (\emph{e.g.} concentration, fraction de phase, teneur en dopant) à partir des spectres/pics, avec \textbf{incertitudes} et \textbf{limites de détection (LOD/LOQ)} robustes. Intégrer à la fois des \textbf{modèles analytiques à base de pics} et des \textbf{régressions multivariées} sous \emph{contraintes physiques}.

\subsection{Notations}
$y(\nu)$ : spectre prétraité (Sec.~5). $S_k$ : aire du pic $k$ (Sec.~6). $C$ : quantité cible (concentration, fraction). $\hat{C}$ : estimateur. $X$ : matrice de \emph{features} (pics, ratios, PCA, \ldots).

\subsection{Quantification univariée (courbe d'étalonnage)}
Pour un pic \emph{marqueur} $k$ d'un analyte $A$ :
\begin{equation}
\mathbb{E}[S_k] \approx \alpha + \beta C_A,\qquad \hat{C}_A = \frac{S_k - \alpha}{\beta}.
\end{equation}
\textbf{Régression pondérée} (WLS) conseillée (hétéroscédasticité). Si saturation: modèle Birks/Hill (logistique) ou spline monotone.
\paragraph{Standard interne} Un pic de référence $r$ stabilise les gains: $R = S_k / S_r \Rightarrow \mathbb{E}[R] \approx \beta' C_A$.
\paragraph{Ratios robustes} Pour plusieurs pics d'un même analyte, utiliser $R_{k} = S_k / \sum_{j\in\mathcal{K}} S_j$ et ajuster $R_k(C)$.

\subsection{LOD/LOQ (ICH Q2)}
Avec le bruit de blanc $\sigma_{\mathrm{blank}}$ et la pente $\beta$ de la courbe d'étalonnage:
\begin{equation}
\mathrm{LOD} = 3.3\, \frac{\sigma_{\mathrm{blank}}}{\beta},\qquad
\mathrm{LOQ} = 10\, \frac{\sigma_{\mathrm{blank}}}{\beta}.
\end{equation}
Alternative : $\sigma$ résiduelle de la régression WLS. Rapporter une IC sur LOD/LOQ par bootstrap.

\subsection{Quantification multipics contrainte}
Soit $\mathbf{S}\in\mathbb{R}^m$ les aires de $m$ pics pour un analyte; on ajuste
\begin{equation}
\mathbf{S} \approx C_A\,\mathbf{h}_A + \boldsymbol{\epsilon},\qquad C_A\ge 0,
\end{equation}
par NNLS (moindres carrés non-négatifs). Avec \textbf{ratios attendus} $R_{A,j}$, ajouter une pénalité $\lambda\sum_j (S_j/\sum S - R_{A,j})^2$.

\subsection{Quantification de phases (mélanges)}
Modèle linéaire de mélange sur \textbf{features} (aires de pics ou spectre réduit):
\begin{align}
\mathbf{y} &\approx \sum_{p=1}^{P} w_p\, \mathbf{y}^{(p)},\quad w_p \ge 0,\ \sum_p w_p = 1,\\
\hat{\mathbf{w}} &= \arg\min_{\mathbf{w}\in\Delta_P}\ \|\mathbf{y} - Y\,\mathbf{w}\|_2^2 + \lambda\|\mathbf{w}\|_2^2,
\end{align}
où $Y=[\mathbf{y}^{(1)},\ldots,\mathbf{y}^{(P)}]$ est la bibliothèque (adaptée à l'instrument, Sec.~8). \textbf{Simplex contraint} via \textsc{Projected Gradient} ou \textsc{CVX}.

\subsection{Régressions multivariées}
\begin{itemize}[nosep]
\item \textbf{PLS} (Partial Least Squares): bon biais-variance, exploite la covariance $X$--$C$.
\item \textbf{ElasticNet}: sélection de variables (pics/PC) sous corrélation; robuste.
\item \textbf{XGBoost/GBRT}: non-linéarités douces; gérer l'interaction entre pics et baseline résiduelle.
\end{itemize}
\textbf{Pondération} des points d'entraînement par l'inverse de la variance de mesure. \textbf{GroupKFold} par lot/site/instrument pour éviter la fuite.

\subsection{Modèles bayésiens hiérarchiques}
\paragraph{Univarié} $S_k \sim \mathcal{N}(\alpha + \beta C, \sigma^2(C))$, priors $\alpha,\beta\sim \mathcal{N}$, $C\ge 0$. Posterior $p(C\mid S_k)$ via MCMC/VI $\Rightarrow$ IC 95\% crédibles.
\paragraph{Mélanges} $\mathbf{y}\sim \mathcal{N}\!\big(\sum_p w_p \mathbf{y}^{(p)},\ \Sigma\big)$, $\mathbf{w}\in\Delta_P$ avec priors Dirichlet. Permet \textbf{incertitudes par phase}.

\subsection{Propagation d'incertitudes}
De l'aire de pic vers $C$ :
\begin{equation}
\mathrm{Var}[\hat{C}] \approx \left(\frac{\partial \hat{C}}{\partial S_k}\right)^2 \mathrm{Var}[S_k] + \left(\frac{\partial \hat{C}}{\partial \beta}\right)^2 \mathrm{Var}[\beta] + \cdots,
\end{equation}
ou par \textbf{bootstrap} (B=500) sur les résidus/échantillons; rapporter IC 95\% et \textbf{couverture} mesurée.

\subsection{Validation croisée \& DOE}
Matrice d'expériences (laser, objectif, temps d'intégration, dopage vrai, \ldots). Split \textbf{par lot/site/instrument}. Métriques: RMSE, MAE, $R^2$, \% within-LOA (Bland--Altman), \textbf{calibration} (pentes proches de 1, intercept proche 0).

\subsection{Sorties normalisées}
\begin{itemize}[nosep]
\item \textbf{quant.csv}: \texttt{sample\_id}, \texttt{analyte} (ou \texttt{phase}), $\hat{C}$, IC95\% $[L,U]$, méthode, LOD/LOQ, flags (OOD, hors domaine).
\item \textbf{calib\_curves.csv}: pentes/intercepts (IC), plage valide, $R^2$, $\sigma$ résiduelle.
\item \textbf{report\_quant.json}: agrégats (RMSE, $R^2$, biais, ECE de probas si bayésien), diagnostics (résidus, leverage).
\end{itemize}

\subsection{Pseudo-code (pipeline quantif)}
\begin{lstlisting}[basicstyle=\ttfamily\small]
def quantify(samples, libs, cfg):
    # 1) Features: pics clés + ratios + éventuelles PC
    X, y_true = build_features(samples, libs, cfg.targets)
    # 2) Choix modèle
    if cfg.model == "PLS":
        mdl = fit_pls(X, y_true, groups=cfg.groups, k=cfg.k)
    elif cfg.model == "ElasticNet":
        mdl = fit_enet(X, y_true, alpha=cfg.alpha, l1=cfg.l1_ratio)
    elif cfg.model == "NNLS_mix":
        mdl = fit_nnls_mix(X, y_true, simplex=cfg.simplex)
    # 3) LOD/LOQ par analyte (courbes)
    lod_loq = compute_lod_loq(mdl, blanks=cfg.blanks)
    # 4) Incertitudes (bootstrap/bayes)
    quant, ci = predict_with_uncertainty(mdl, X, method=cfg.uncertainty)
    # 5) Exports
    write_csv("artifacts/quant.csv", quant)
    write_csv("artifacts/calib_curves.csv", summarize_calib(mdl))
    write_json("artifacts/report_quant.json", make_report(quant, ci, lod_loq))
    return quant
\end{lstlisting}

\subsection{Tests \& DoD}
\begin{itemize}[nosep]
\item \textbf{Exactitude}: RMSE $\downarrow$, $R^2\uparrow$; \textbf{Biais} $|\hat{C}-C|$ médian proche 0.
\item \textbf{Robustesse}: invariance au gain/offset (avec standard interne), tolérance au bruit $\pm$10--20\%.
\item \textbf{LOD/LOQ}: cohérents avec la variance de blanc; IC bootstrap étroites et stables.
\item \textbf{Couverture}: IC 95\% couvrent $\ge 93$--$97$\% des vérités sur val. indépendante.
\item \textbf{DoD}: fichiers \texttt{quant.csv}, \texttt{calib\_curves.csv}, \texttt{report\_quant.json} produits et versionnés.
\end{itemize}

\subsection{Risques \& parades}
\begin{itemize}[nosep]
\item \textbf{Effets de matrice}: privilégier ratios/standard interne; inclure covariables (tint, laser) dans $X$.
\item \textbf{Non-linéarité}: modèles splines/GBRT; limiter la plage de validité et l'annoncer.
\item \textbf{Déséquilibre}: stratifier les splits, pondérer les pertes; data augmentation physiquement plausible.
\item \textbf{Dérive instrumentale}: recalibration régulière et harmonisation (Sec.~4); AL pour ré-étalonner.
\end{itemize}

% Table des colonnes normalisées pour quant.csv

\begin{table}[h]
\centering
\caption{Colonnes normalisées du fichier \texttt{quant.csv}.}
\begin{tabular}{ll}
\toprule
\textbf{Colonne} & \textbf{Description} \\
\midrule
sample\_id & Identifiant d'échantillon \\
analyte & Nom de l'analyte ou de la phase \\
C\_hat & Estimation de la quantité (unités SI ou fraction) \\
CI95\_L & Borne basse de l'intervalle de confiance 95\% \\
CI95\_U & Borne haute de l'intervalle de confiance 95\% \\
method & Méthode employée (PLS, ENet, NNLS\_mix, Bayes, \ldots) \\
LOD & Limite de détection (mêmes unités que C) \\
LOQ & Limite de quantification \\
flags & Indicateurs (OOD, hors domaine, saturation, \ldots) \\
meta & Champs utiles (standard interne, ratio utilisé, version modèle) \\
\bottomrule
\end{tabular}
\end{table}

